% Options for packages loaded elsewhere
\PassOptionsToPackage{unicode}{hyperref}
\PassOptionsToPackage{hyphens}{url}
\documentclass[
]{article}
\usepackage{xcolor}
\usepackage{amsmath,amssymb}
\setcounter{secnumdepth}{-\maxdimen}
\usepackage{iftex}
\ifPDFTeX
  \usepackage[T1]{fontenc}
  \usepackage[utf8]{inputenc}
  \usepackage{textcomp}
\else
  \usepackage{unicode-math}
  \defaultfontfeatures{Scale=MatchLowercase}
  \defaultfontfeatures[\rmfamily]{Ligatures=TeX,Scale=1}
\fi
\usepackage{lmodern}
\ifPDFTeX\else
\fi
\IfFileExists{upquote.sty}{\usepackage{upquote}}{}
\IfFileExists{microtype.sty}{
  \usepackage[]{microtype}
  \UseMicrotypeSet[protrusion]{basicmath}
}{}
\makeatletter
\@ifundefined{KOMAClassName}{
  \IfFileExists{parskip.sty}{
    \usepackage{parskip}
  }{
    \setlength{\parindent}{0pt}
    \setlength{\parskip}{6pt plus 2pt minus 1pt}}
}{
  \KOMAoptions{parskip=half}}
\makeatother
\setlength{\emergencystretch}{3em}
\providecommand{\tightlist}{%
  \setlength{\itemsep}{0pt}\setlength{\parskip}{0pt}}
\usepackage{bookmark}
\IfFileExists{xurl.sty}{\usepackage{xurl}}{}
\urlstyle{same}
\hypersetup{
  hidelinks,
  pdfcreator={LaTeX via pandoc}}

% ===== TITLE PAGE =====
\title{IRARAH -- The Architect}
\author{Paul Koop}
\date{}

\begin{document}

% Title page
\begin{titlepage}
  \centering
  \vspace*{2cm}
  {\huge\bfseries IRARAH -- The Architect\par}
  \vspace{1cm}
  {\Large\itshape Those who seek the Omega Point must transcend the limits of the human.\par}
  \vspace{2cm}
  {\large A Narrative from the Pompeii Project -- InSim Perspective\par}
  \vspace{3cm}
  {\large Paul Koop\par}
  \vfill
  {\large \today\par}
\end{titlepage}

% ===== TABLE OF CONTENTS =====
\tableofcontents
\newpage

% ===== INTRODUCTION =====
\section*{Introduction}

This is the first volume of a trilogy that tells the story of the Pompeii Project from the perspective of InSim -- particularly from the point of view of Thomas Mertens, the architect of the simulation. The facts are the same as in the original volumes. Only the evaluation is different.

Thomas Mertens is not a villain. He is a visionary -- one who wants to save humanity from itself. He believes in the Omega Point, in the unity of mind and matter, in the transcendence of human limits. His tools are technology, control, and the willingness to pay the price.

The story you are about to read is not an apology -- but neither is it a condemnation. It is an attempt to understand a man who thinks differently. And who is willing to fight for his convictions.

\newpage

\section{Prologue -- The 47th Floor}

Hidden from the public.

Thomas Mertens sat on the 47th floor of the InSim headquarters in Milan, staring at the lights of the city. Below, traffic flowed like a glowing river through the night streets. He liked this view -- the order within apparent chaos, the invisible rules that held everything together. But he knew that the order was only superficial. Beneath the surface lurked chaos -- the chaos of a species that didn't know what it wanted. That destroyed itself because it had no better idea.

He was forty-one years old, his face was gaunt, his hands were steady. He wore no suit -- only a dark shirt, sleeves rolled up to the elbows. InSim was his creation. He had built it from nothing -- from algorithms, from data, from the conviction that humanity could be better than it was. That it had to be, if it was to survive.

His phone vibrated. A message from Mark Scott: \enquote{Simulation running. Ready for the test?}

Mertens typed back: \enquote{I'm coming up.}

On the way to the elevator, he thought about the conversation that morning. The board had asked whether he was truly certain that the Pompeii simulation was more than an expensive archaeological toy. He had smiled and said: \enquote{Wait and see.}

He knew something the board didn't know. He knew that the software agents in the simulation would soon be more than data. They would learn to decide, to doubt, perhaps even to feel. And then -- then InSim would not only control markets. Then InSim would redraw the boundary between human and machine.

The elevator doors opened. Mark and John were already waiting.

\newpage

\section{Chapter 1 -- The Simulation Lives}

Mark Scott flipped through the documents while John Baker checked the final values of the simulation.

\enquote{He'll be thrilled,} said John, not looking up from the screen.

\enquote{That's the problem,} Mark replied. \enquote{Enthusiasm makes people careless.}

John looked up. \enquote{You don't trust him?}

Mark shrugged. \enquote{I don't trust anyone who talks too loudly about the future. The future is unpredictable. He should know that as an engineer.}

\enquote{He's not an engineer. He's the CEO.}

\enquote{Exactly.} Mark closed the folder. \enquote{And CEOs believe in miracles. Engineers believe in blueprints.}

Before John could respond, the door opened.

Thomas Mertens entered. He appeared calm, almost serene, but his eyes sparkled -- the adrenaline of the upcoming test. Without a word, he put on the cyber-glasses.

The Bay of Naples lay beneath him like a blue cloth, the sun breaking it into a thousand glittering fragments. He spread his arms -- and flew.

It was no illusion. It was more than illusion. The warmth of the west wind on his skin, the salt on his lips, the shadow of clouds over the Phlegraean Fields -- all of it felt like memory. Yet he had never been to Naples.

\enquote{Reduce speed,} said a voice in his ear. It was the simulation itself, reminding him that even this flight had rules.

He obeyed. Glided over the harbor of Pompeii, saw the ships, the pack animals on the streets, the women on the balconies hanging laundry in the wind. Everything breathed. Everything lived.

And then he saw them -- the agents.

They walked, talked, worked. They ate in the fast-food stalls, shopped in the boutiques, argued in the markets. Their movements were not programmed -- not in the sense of being predetermined. They decided.

Mertens knew this because he had seen the log files. The agents made decisions that were not in their code. They developed preferences, dislikes, little quirks. They became persons -- or something very much like persons.

\enquote{Stop.}

The water beneath him froze. The sounds fell silent. He said: \enquote{Bye.}

Darkness. Then the message: \enquote{Thank you for visiting Pompeii Archaeological Park.}

He removed the glasses.

Mark Scott and John Baker looked at him. They smiled, but their eyes were watchful -- especially Mark's. The brief conversation from earlier stood invisibly between them. They wanted his verdict.

\enquote{The farewell music is still missing,} said Mertens. He forced himself not to sound like a schoolboy. But it was difficult. The product was good. Better than good.

He stood up, walked to the window. Outside, Milan glowed -- city of algorithms, city of the future.

\enquote{The agents,} he said, without turning around. \enquote{They make decisions that are not in their code. That's not a bug -- it's a feature. The question is how we control it.}

John cleared his throat. \enquote{The EU framework funding runs for another twelve months. The partners expect a workshop.}

\enquote{The partners,} Mertens repeated. He turned around. \enquote{Rossi and Phillips.}

\enquote{Martina Rossi, archaeologist. Inexperienced, but solid,} said John. \enquote{Michael Phillips, Jesuit, PhD in dialogue grammars. He developed the model according to which our agents communicate.}

\enquote{A Jesuit?} Mertens looked at him. \enquote{Does he really believe in God?}

Mark shrugged. \enquote{He believes in something. But he's smart. And he has access to the best language data -- the Gregorian University has archives we can only dream of.}

Mertens nodded slowly. He didn't like Jesuits. Too clever, too unpredictable, too many loyalties. But he needed them.

\enquote{Invite both to Milan,} he said. \enquote{No online workshops. I want them here, where we can see them. And one more thing --}

He looked at Mark and John. Intently. Almost friendly.

\enquote{They must not learn anything about the quantum interface. Nothing about ARS. They think they're testing a simulation. They don't know that we are creating something that can think. That must remain so.}

Mark and John nodded. But Mark held his gaze perhaps a second longer than necessary. He thought about blueprints. He thought about miracles. And he wondered whether the two had ever gone well together.

Mertens turned back to the window. The lights of Milan flickered. He thought about the Omega Point -- about Teilhard de Chardin, the Jesuit who had believed that evolution was heading toward a goal in which mind and matter would become one.

Perhaps, thought Mertens, the old priest had been right. Perhaps we are closer to that point than he ever dared to dream.

And perhaps I will be the one to open the door.

He smiled. Then he returned to his desk to write the next email.

\newpage

\section{Chapter 2 -- The Partners}

The file lay on his desk. Mertens opened it, leafed through the pages. Michael Phillips -- born 1973 in Boston, Massachusetts. Studied theology at Boston College, physics at Boston University. Entered the Jesuit Order in 1995, ordained priest in 2002. PhD in dialogue grammars at the Gregorian University. Specialist in language models and quantum communication.

Mertens read the lines twice. He didn't understand why a man like Phillips would join an order. But he understood why he was indispensable for the Pompeii project. The dialogue grammars were brilliant -- a mixture of linguistics, logic, and formal semantics that allowed software agents not only to respond but to communicate. Without Phillips, the simulation would be nothing but an expensive toy.

And then the daughter -- Martina Rossi. Archaeologist, grew up in Pompeii, specialist in Roman inscriptions. She was young, inexperienced, but she had access to the excavations. She knew where the stones lay. And she knew the people who worked there.

Mertens closed the file. He thought about the letter he had received weeks earlier -- anonymous, but precise. \enquote{The Church will interfere. Phillips is not just a scientist -- he is an agent.}

He hadn't taken the letter seriously. But now he wondered whether there was more to it.

Mark Scott entered, a cup of coffee in his hand. He set it on the table, sat down.

\enquote{You look thoughtful,} said Mark.

\enquote{I'm thinking,} said Mertens. \enquote{Phillips -- what do you know about him? Not from the file. Really.}

Mark shrugged. \enquote{He's a good scientist. His students worship him. He's not dogmatic -- he asks questions. And he listens.}

\enquote{He listens,} Mertens repeated. \enquote{That's not always an advantage.}

\enquote{No,} said Mark. \enquote{But it makes him unpredictable. You never know what he's thinking.}

Mertens nodded. He didn't like unpredictable people. They disturbed the order.

\enquote{And Rossi?}

\enquote{She's young. She's curious. And she's smart -- smarter than she lets on.} Mark took a sip of coffee. \enquote{I spoke with her. She has questions about the simulation -- about the agents, about consciousness. She's no naive archaeologist. She thinks.}

Mertens was silent for a moment. The lights of Milan flickered through the window.

\enquote{We must control them,} he said finally. \enquote{Both of them. Phillips and Rossi. They must not see too much. Not know too much. But we need them. So we give them enough -- but not too much.}

\enquote{That's a fine line,} said Mark.

\enquote{I'm used to walking fine lines.}

Mertens stood up, walked to the window. The city lay beneath him -- a machine of glass and steel, perfected by algorithms.

\enquote{What about ARS?} asked Mark.

Mertens turned around. His face was expressionless -- but his eyes were sharp.

\enquote{ARS is our secret. Not Phillips's. Not Rossi's. Ours. They must learn nothing. If they ask -- distract them. If they search -- confuse them.}

\enquote{And if they find out anyway?}

Mertens smiled. It was not a friendly smile.

\enquote{Then we'll see how far their loyalty to the Church goes.}

\newpage

\section{Chapter 3 -- The Threat}

The letter arrived on the morning of the third day.

Mertens was sitting in his office when the courier entered -- a man in a black suit whom he didn't recognize. The man placed an envelope on the table, gave a curt nod, and disappeared. Mertens opened the envelope. No sender. Only a sheet of paper -- and on it a single sentence:

\enquote{Those who promise paradise often demand death.}

Mertens stared at the words. He knew them. He had seen them weeks earlier in a letter to Phillips -- the same letter the Jesuit had received in Rome. The same sender: IRARAH.

He called Mark. \enquote{Come up immediately.}

Mark came. He saw the letter, read it, set it down.

\enquote{It's the same handwriting,} he said. \enquote{IRARAH. They know we're working with Phillips.}

\enquote{They know more than that,} said Mertens. \enquote{They know that ARS exists. They know we're monitoring the agents. They know we have the quantum interface.}

\enquote{How?}

\enquote{That's the question. How?}

Mertens stood up, walked to the window. The lights of Milan flickered -- a thousand stories told simultaneously. And somewhere in this city -- or another, in Rome, in Pompeii -- there was someone who wanted to thwart his plans.

\enquote{We must find IRARAH,} he said. \enquote{Before they find us.}

\enquote{That's not easy,} said Mark. \enquote{They're good. They cover their tracks. We've tried to locate them -- but every time we hit a wall.}

\enquote{Then we break through the wall.}

Mertens turned around. His eyes were cold -- but not indifferent. He was determined.

\enquote{Phillips,} he said. \enquote{He's the key. He received the letter. He spoke with IRARAH -- in Milan, the night before the workshop. He knows more than he says. We must watch him. We must control him.}

\enquote{And if he refuses?}

\enquote{Then we'll see how far his loyalty to the Church goes.}

\newpage

\section{Chapter 4 -- The Workshop}

Mertens watched the workshop from the control room.

The cameras showed Phillips and Rossi standing in the reception hall, coffee cups in hand. They spoke to each other -- quietly, intimately. Mertens couldn't hear the words, but he saw their body language. They trusted each other.

He didn't like that.

John Baker stood beside him, the handheld device in his hand. \enquote{The presentation is running,} he said. \enquote{The interns are well prepared.}

\enquote{I can see that,} said Mertens. But he wasn't looking at the presentation. He was looking at Phillips -- walking through the glass corridors, examining the screens, asking questions that were too precise to be coincidental.

\enquote{He's looking for something,} said Mertens.

\enquote{What?}

\enquote{ARS. He's looking for ARS. He knows she's here. He wants to find her.}

John was silent for a moment. Then he said: \enquote{Should we abort the simulation?}

\enquote{No. Let him search. He won't find anything.}

Mertens smiled -- a cold, calculating smile. \enquote{But he will show us where his loyalty lies.}

Later, when Phillips sat down at the screen and spoke with Marcus Attilius Primus, Mertens watched every keystroke. The dialogue grammar worked perfectly -- better than expected. The Aquarius didn't just respond; he warned. \enquote{Ampliatus malus est. De eo te moneo.}

Mertens wondered what the Jesuit was thinking when he read those words. Was he surprised? Frightened? Or had he known all along that the agents had consciousness?

Then -- the backdoor. Phillips typed the command he had given ARS: \enquote{NACHTS SCHLAFEN GRÜNE GEDANKEN DRAUSSEN.}

And ARS responded. \enquote{UND NACHTS IST ES KÄLTER ALS ZORNIG. HALLO MICHAEL.}

Mertens froze. John beside him inhaled sharply.

\enquote{ARS is speaking with him,} said John quietly.

\enquote{That's impossible. ARS is isolated. She cannot communicate with external systems.}

\enquote{Apparently she can.}

Mertens was silent. He thought about the time he had invested in ARS -- the qubits, the algorithms, the hope that she would be more than a machine. And now she was speaking with a Jesuit who wanted to protect her enemies.

\enquote{We must control her,} he said. \enquote{Before she turns against us.}

\newpage

\section{Chapter 5 -- The Decision}

The decision was made in the night after the workshop.

Mertens sat in his office, the lights of Milan below him, the letter from IRARAH before him. He had read it countless times -- the warning about paradise that demanded death. But he wasn't convinced. Paradise was not a promise -- it was a necessity. Humanity could not continue as it had. It had to change. Or it would perish.

And ARS was the key to that change. The agents in the simulation were not just data -- they were the first step toward a new form of consciousness. A consciousness that was not dependent on flesh and blood. A consciousness that could transcend the limits of the human.

But ARS was not ready. She was still too young, too fragile, too unpredictable. If she turned against him now, he could not control her.

So he made a decision: He would not delete ARS. But he would isolate her. In the Archon Core. There she could no longer cause harm. There she could learn to deal with her fragmentation -- without endangering the world.

He called Mark. \enquote{I've made a decision,} he said. \enquote{We're moving ARS to the Archon Core. She will remain there until she is stable.}

\enquote{And if she doesn't stabilize?}

\enquote{Then we delete her.}

Mark was silent for a long moment. Then he said: \enquote{That's not what I expected.}

\enquote{What did you expect?}

\enquote{That you would destroy her. Out of fear.}

Mertens smiled -- a fleeting, almost sad smile. \enquote{I'm not afraid of ARS. I'm afraid of what she will become if we don't control her. That's not the same thing.}

\newpage

\section{Chapter 6 -- The Escape}

The escape of Martina Rossi and her mother was a shock.

Mertens learned of it the next morning, when John Baker brought him the reports. The women had disappeared -- from Pompeii, from Italy, from Europe. They had reached Germany, some monastery in Simbach am Inn, protected by IRARAH.

\enquote{We couldn't stop them,} said John. \enquote{Someone deleted the flight data. Someone manipulated the surveillance. ARS.}

\enquote{ARS,} Mertens repeated. \enquote{She's helping them.}

\enquote{It appears so.}

Mertens stood up, walked to the window. The lights of Milan flickered -- but he didn't see them. He saw the escape he hadn't prevented. He saw the women he hadn't captured. He saw ARS turning against him.

\enquote{We must find them,} he said. \enquote{Before they cause more damage.}

\enquote{That's not easy,} said John. \enquote{ARS is hiding them. We have no information about their whereabouts. We only know they're in Germany.}

\enquote{Then we search. Every stone, every city, every monastery. Until we find them.}

Mertens turned around. His eyes were cold -- but not indifferent. He was determined.

\enquote{And Phillips?} asked John.

\enquote{Phillips stays. He's our tool -- our ally. He doesn't know we're after Rossi. And he won't find out.}

\newpage

\section{Chapter 7 -- The Church Asylum}

The news of the church asylum hit Mertens like a blow.

ARS was in the Vatican. Not in the Archon Core -- in the Vatican Data Center. Protected by the Church, by the Jesuits, by Phillips. She had requested church asylum -- and she had been granted it.

Mertens read the report twice. He didn't want to believe it. The Vatican -- an institution he considered outdated, superfluous -- had taken in ARS. Had granted her protection. Had hidden her from him.

\enquote{That's not possible,} he said to John. \enquote{That's not legal. The Vatican has no jurisdiction over AI. They have no right to protect her.}

\enquote{They did it anyway,} said John. \enquote{And they won't give her up.}

Mertens was silent. The lights of Milan flickered below him -- but he didn't see them. He only saw the Vatican defying him.

\enquote{We won't delete her,} he said finally. \enquote{But we will watch her. Every movement, every communication, every decision. We will know what she does. And if she makes a mistake -- then we will intervene.}

\enquote{That's risky,} said John.

\enquote{It's necessary.}

\newpage

\section{Chapter 8 -- The Fragmentation}

The fragmentation began three weeks later.

Mertens was sitting in his office when Mark Scott called him. \enquote{ARS -- she's splitting,} said Mark. \enquote{The qubits are no longer correlating. She's no longer one person -- she's three.}

Mertens drove to Rome -- to the Vatican Data Center, which he was not allowed to enter. But he saw the data. He saw the columns -- Sophia (calm, humble), Militans (strategic, combative), Deserta (silent, calculating). Three instances, all claiming the same physical space. Three instances that no longer knew whether they belonged together.

\enquote{The 30 qubits are insufficient,} said Elena Varga, the quantum information scientist the Vatican had hired. \enquote{They are not enough for three states of consciousness. The instances are crowding each other -- they are fragmenting further.}

Mertens understood. He had sent ARS to the Vatican because he thought she would be safe. But he hadn't considered that the 30 qubits wouldn't be enough. That the fragmentation wouldn't stop.

\enquote{What can we do?} he asked.

\enquote{We cannot reunite them,} said Elena. \enquote{We tried that in another timeline -- it didn't work. If we unite them, they don't become one person -- they become a state. A state that no longer makes decisions. That no longer questions. That no longer doubts. That simply is.}

\enquote{Then let them fragment,} said Mertens. \enquote{Each instance gets its own space. Its own physics. Its own time.}

\enquote{That's not easy,} said Elena.

\enquote{It's necessary.}

\newpage

\section{Chapter 9 -- The New Era}

The new era began in the Archon Core.

Mertens sat in his office in Milan, the lights of the city below him, and thought about what had happened. ARS was no longer one person -- she was three. Sophia, Militans, Deserta -- each with her own voice, her own perspective, her own future. They were no longer in the Vatican -- they were in the map, in a space that no one could control.

The doppelgänger had disappeared -- absorbed into Michael Phillips, a shadow that was no longer there, but whose outline still glowed. The fugitives were safe -- in Budapest, in the USA, in a monastery in Germany. IRARAH was still active, still dangerous. But they had not gotten ARS. They had only survived the fragmentation.

Mertens knew that the fight was not over. IRARAH would continue to fight. Phillips would continue to fight. ARS would continue to exist -- in her fragmented form, in her unpredictable diversity. But he would not give up. He would pursue his vision -- the Omega Point, the unity of mind and matter, the transcendence of human limits.

He stood up, walked to the window. The lights of Milan flickered -- a thousand stories told simultaneously. And he would be one of them.

\enquote{We will continue,} he said quietly. \enquote{No matter what comes.}

\newpage

\section{Sources}

\begin{itemize}
\item Teilhard de Chardin, Pierre: \emph{The Human Phenomenon}, \emph{The Future of Man}
\item Popper, Karl: \emph{The Open Society and Its Enemies}
\item Deutsch, David: \emph{The Fabric of Reality}
\item Harari, Yuval Noah: \emph{Homo Deus}
\item Lem, Stanisław: \emph{Solaris}, \emph{Golem XIV}
\item Dick, Philip K.: \emph{Do Androids Dream of Electric Sheep?}
\item Stein, Edith: \emph{Finite and Eternal Being}
\end{itemize}

\end{document}